% Fase 1 --- Propósito, requisitante, partes interessadas e escopo (responsável: Heitor).
% Versão completa para a GitPage: docs/fase1-proposito.md.

\subsubsection{Passo 1.1: Propósito, Requisitante e Partes Interessadas}
\label{subsec:proposito}
\label{subsec:stakeholders}

O Detran-DF oferece a consulta de infrações no Portal de Serviços e no aplicativo móvel. Embora compartilhem o servidor de identidade, os canais consultam serviços de \textit{backend} distintos (Seção~\ref{subsec:produto}), o que pode gerar diferenças de completude, tempo de resposta ou conclusão da tarefa. Esta avaliação comparativa tem como propósito verificar se os dois canais entregam qualidade equivalente e apoiar quatro decisões institucionais da \textbf{Área de Governança Digital do Detran-DF} (requisitante hipotético adotado para orientar a avaliação, sem vínculo oficial com o órgão), sintetizadas na Tabela~\ref{tab:decisoes}.

\begin{table}[htbp]
\centering
\caption{Decisões apoiadas pela avaliação e seus responsáveis}
\label{tab:decisoes}
\footnotesize
\setlength{\tabcolsep}{3pt}
\begin{tabular}{|c|p{6.5cm}|p{3.4cm}|p{2.7cm}|}
\hline
\textbf{ID} & \textbf{Decisão} & \textbf{Quem decide} & \textbf{Evidência usada} \\ \hline
D1 & Em qual canal e característica investir primeiro & Requisitante & Objetivos G1, G2 e G3 \\ \hline
D2 & Quais divergências de dados corrigir com urgência & Equipe técnica de TI & Questão Q1.2 \\ \hline
D3 & Qual canal recomendar ao cidadão em cada contexto & Requisitante e Atendimento & Objetivos G2 e G3 \\ \hline
D4 & Quais plataformas e redes devem ter suporte garantido & Equipe técnica de TI & Objetivo G3 \\ \hline
\end{tabular}
\par\smallskip
\raggedright\scriptsize Fonte: elaborada pela equipe.
\end{table}

A Tabela~\ref{tab:stakeholders} sintetiza as partes interessadas mapeadas, seus critérios de sucesso e sua influência direta na avaliação (o detalhamento estendido encontra-se na GitPage em \texttt{docs/fase1-proposito.md}). As três primeiras fornecem os pontos de vista dos objetivos GQM (Seção~\ref{subsec:rastreabilidade}).

\begin{table}[htbp]
\centering
\caption{Síntese das partes interessadas, critérios de sucesso e influência}
\label{tab:stakeholders}
\footnotesize
\setlength{\tabcolsep}{2.5pt}
\begin{tabular}{|p{3.0cm}|p{3.4cm}|p{3.6cm}|p{3.2cm}|}
\hline
\textbf{Parte interessada} & \textbf{Necessidade principal} & \textbf{Critério de sucesso} & \textbf{Influência (GQM)} \\ \hline
Requisitante: Gov. Digital Detran-DF (hipotético) & Verificar paridade entre canais e priorizar melhorias & Julgamento por canal para apoiar D1 & Propósito comparativo transversal \\ \hline
Cidadão/condutor proprietário do veículo & Dados completos e corretos para pagar ou recorrer & Consulta concluída e consistente nos dois canais & Adequação funcional e contexto (G1 e G3) \\ \hline
Equipe técnica de TI do Detran-DF & Identificar lentidão, consumo de rede ou falhas & Medições reproduzíveis por canal e rede (D2, D4) & Eficiência de desempenho (G2) \\ \hline
Atendimento, Controle e Equipe Avaliadora & Orientar uso (D3), auditar transparência e viabilizar coleta & Recomendação por contexto e execução em caixa-preta & Comparação em G2/G3 e cortes de escopo \\ \hline
\end{tabular}
\par\smallskip
\raggedright\scriptsize Fonte: elaborada pela equipe. Tabela completa na GitPage (\texttt{docs/fase1-proposito.md}).
\end{table}

\subsubsection{Passo 1.2: Escopo, Profundidade e Limites de Validade}
\label{subsec:escopo}

O objeto da avaliação é o \textbf{histórico de infrações por veículo}, na área autenticada, no Portal de Serviços \cite{detrandf2026portal} e no aplicativo (v.~2.6.9, Android e iOS) \cite{detrandf2026app}, observado da escolha do menu de infrações (pós-login) até a exibição dos detalhes da autuação. Os contextos avaliados compreendem: (i) \textbf{Portal:} Chrome e Firefox em desktop, e Chrome em Android; (ii) \textbf{Aplicativo:} smartphones Android de entrada e intermediário (Android $\ge 8.0$) e iPhone (iOS $\ge 15.1$); e (iii) \textbf{Rede:} Wi-Fi, 4G e rede degradada simulada.

A avaliação adota profundidade de caixa-preta com métricas estritamente externas. A Tabela~\ref{tab:fora-escopo} sintetiza os cortes de escopo. Como limites de validade, a massa de dados restringe-se a veículos reais com consentimento escrito de integrantes e familiares, os tempos refletem o período de coleta sem controle de carga dos servidores, e a correção funcional baseia-se na consistência mútua entre os canais.

\begin{table}[htbp]
\centering
\caption{Itens fora do escopo e justificativa sintética de cada corte}
\label{tab:fora-escopo}
\footnotesize
\setlength{\tabcolsep}{3pt}
\begin{tabular}{|p{4.2cm}|p{8.8cm}|}
\hline
\textbf{Item fora do escopo} & \textbf{Justificativa do corte} \\ \hline
Autenticação (login) e Consulta pública sem login & Login usa mecanismos incomparáveis (reCAPTCHA vs. biometria); consulta pública omite o detalhamento da autuação. \\ \hline
Histórico por condutor e Veículos de outra UF & Consulta por condutor existe apenas no aplicativo; outra UF depende de bases externas ao Detran-DF. \\ \hline
Pagamento, recursos e indicação de condutor & Constituem serviços transacionais distintos que exigiriam operações financeiras ou administrativas reais. \\ \hline
Capacidade de interação (usabilidade) e SOs antigos & Usabilidade excluída pelas diretrizes da disciplina; Android $< 8.0$ e iOS $< 15.1$ não são suportados pela loja. \\ \hline
\end{tabular}
\par\smallskip
\raggedright\scriptsize Fonte: elaborada pela equipe. Detalhamento completo em \texttt{docs/fase1-proposito.md}.
\end{table}